\documentclass[11pt,a4paper]{article}
\usepackage[utf8]{inputenc}
\usepackage[T1]{fontenc}
\usepackage{hyperref}
\usepackage{booktabs}
\usepackage{graphicx}
\usepackage{xcolor}
\usepackage[margin=1in]{geometry}
\hypersetup{colorlinks=true, linkcolor=blue, urlcolor=blue, citecolor=blue}
\title{Telugu-First TokEval: Small Dedicated Tokenizers Beat Massive Multilingual Vocabularies\\ on Fertility and Parity}
\author{Gottipati Jamadagni\\ Software Engineer, Unnanu, Hyderabad, India\\ \texttt{jama@unnanu.com} \and AI Collaborator (Muse Spark)}
\date{August 23, 2026}
\begin{document}
\maketitle
\begin{abstract}
Byte-Pair Encoding tokenizers optimized for English incur an $11.9\times$ fertility penalty on Telugu: \texttt{tiktoken cl100k} averages 13.57 tokens/word on Telugu vs 1.14 on English. Massive multilingual vocabularies (Sarvam 262K, Saiteja 50K) partially close the gap to 2.48 and 1.64. We replicate the TokEval suite (Meister et al., COLM 2026) on Telugu using 80K Telugu Wikipedia lines and FLORES-200 997-sentence parallel set. Dedicated SentencePiece BPE/Unigram at 16K/32K achieve \textbf{1.55 tokens/word} --- $8.8\times$ longer effective context than tiktoken and 37\% better than Sarvam 262K despite $8\times$ smaller vocabulary. Unigram preserves morphology better. Joint Telugu-English 16K balances parity to 1.34 with 64 MB embedding vs 2.1 GB, fitting phone-scale LLMs. Code and tokenizers released.
\end{abstract}
\noindent\textbf{Keywords:} tokenization, Telugu, TokEval, parity, edge LLM

\section{Introduction}
Phone-scale 1B--2B LLMs that learn collaboratively need Telugu tokenization that does not burn 6--12 tokens per word. Recent TokEval \cite{tokeval} and Indic tokenization studies \cite{brahma2026} give metrics, but no Telugu replication at Wikipedia scale with phone-fit analysis exists. We fill that gap.

\section{Related Work}
TokEval \cite{tokeval}, Brahma et al. \cite{brahma2026}, Parity-aware BPE \cite{foroutan2026,kanjirangat2026}, Sarvam \cite{sarvam}, FreeToken \cite{freetoken}.

\section{Data and Method}
Corpora: \texttt{wikimedia/wikipedia 20231101.te} 80K lines (36 MB), 20K English, FLORES-200 997 parallel. Method: SentencePiece BPE/Unigram 16K/32K, joint vs dedicated, TokEval metrics: fertility, compression, parity, digit handling.

\section{Results}
\begin{table}[h]
\centering
\small
\begin{tabular}{lrrrrrl}
\toprule
Name & Vocab & TE & EN & Parity & TE-wiki & Embed\\
\midrule
te\_bpe\_32k & 32000 & \textbf{1.55} & 2.47 & 0.62 & 1.85 & 128 MB\\
te\_uni\_32k & 32000 & \textbf{1.55} & 2.51 & 0.62 & 1.85 & 128 MB\\
Saiteja 50K & 50000 & 1.64 & 1.28 & 1.28 & 1.47 & 200 MB\\
te\_uni\_16k & 16000 & 1.68 & 2.96 & 0.57 & 2.07 & 64 MB\\
joint\_bpe\_16k & 16000 & 1.91 & 1.42 & 1.34 & 2.29 & 64 MB\\
Sarvam 262K & 262144 & 2.48 & 1.19 & 2.08 & 2.79 & 2147 MB\\
tiktoken & 100256 & 13.57 & 1.14 & 11.90 & 13.89 & 400 MB\\
\bottomrule
\end{tabular}
\caption{Telugu fertility (tokens/word, lower better). Embedding = vocab$\times$2048$\times$2B.}
\end{table}

Insights: $8.8\times$ win over tiktoken, dedicated $>$ massive multilingual, Unigram morphology, joint tradeoff, digit split failure.

\section{Analysis}
Phone fit: 64 MB vs 2.1 GB enables 1.7B Q4\_KM on Android. TokEval maps to 8.7$\times$ bits-per-byte gain predicted.

\section{Limitations}
No pretraining proxy yet, no Tamil/Kannada cluster, digit handling.

\section{Release}
Code: \texttt{exp\_telugu\_tokeval/} MIT, HF Hub pending.

\begin{thebibliography}{9}
\bibitem{tokeval} Meister et al. TokEval. \url{https://arxiv.org/abs/2608.18062}
\bibitem{brahma2026} Brahma et al. Findings ACL 2026. \url{https://aclanthology.org/2026.findings-acl.1632/}
\bibitem{kanjirangat2026} Kanjirangat et al. MeLLM 2026. \url{https://aclanthology.org/2026.mellm-1.21/}
\bibitem{sarvam} Sarvam 30B. \url{https://huggingface.co/sarvamai/sarvam-30b}
\bibitem{freetoken} FreeToken. \url{https://arxiv.org/abs/2608.16157}
\end{thebibliography}
\end{document}
